\subsection{PRISM}\label{sec:2.2}

% Sources: SYSTEM_PLAN.md section 1.
PRISM turns one follow-up scan and the patient's clinical record into a structured record and a checked report in
six stages (Figure~\ref{fig:1}). It outlines the tumour, measures it, scores the scan with three models, attaches
statements of how certain each score is, collates every value into one record with its source, and has a local
language model write the report, which rule-based programs then check against the record. Table~\ref{tab:models}
lists every score in the paper.

% Sources: OUTLINE.md (Fig 1); SYSTEM_PLAN.md section 1.
% FIGURE REDESIGN (2026-11-18), modelled on Tak et al. 2025, NEJM AI, doi:10.1056/AIoa2400703.
% Example paper (local only): docs/literature/examples/tak2025_nejmai.md and tak2025_figures/.
% What we take from it, figure by figure: docs/literature/examples/tak2025_NOTES.md.
% Figure 1a follows their Figure 1A (one real patient on a timeline). Drawn by scripts/fig1_prism.py, which reads the
% MU images from data/mu-glioma-post (patient PatientID_0204; labels by the paper's rule, computed in the script).
\begin{figure}[!htbp]
\centering
\figslot{fig1_system.pdf}{0.6\textheight}{One MU patient's follow-up on real scans; PRISM's segmenter; the method in six stages; the evaluation protocol.}
\caption{\textbf{The task and the method.} (a) One MU-Glioma-Post patient's scans after radiotherapy. Each panel is
the T1c axial slice with the most enhancement. The enhancing tumour is outlined as PRISM's segmenter, a Swin UNETR
trained on BraTS 2021 \citep{hatamizadeh2022}, draws it (solid) and as the mask supplied with the dataset draws it
(dashed); volumes are the segmenter's, with the supplied mask's in brackets. Enhancement at 3.7
weeks, inside the first 12 weeks, receded without progression; progression was documented at 33.4 weeks. Scans are
labelled by the paper's rule: progression documented at the scan, or no progression for at least 120 days after it;
the scan at 24.7 weeks meets neither and is not labelled. (b) PRISM's segmenter, shown on the patient's last scan. A
Swin UNETR \citep{hatamizadeh2022} takes the four sequences as channels. Its Swin Transformer encoder applies
self-attention within shifted local windows and halves the resolution at each stage, and its convolutional decoder
restores it through skip connections from every stage. Each Swin stage holds two successive blocks (below):
layer normalisation, multi-head self-attention within 7\textsuperscript{3} windows, and a multilayer perceptron, each
with a residual connection; the second block shifts the windows by half a window, so that information passes between
neighbouring windows. Three sigmoid outputs (tumour core, whole tumour, enhancing tumour) give the
enhancing, necrosis and FLAIR masks. (c) What PRISM computes at one follow-up scan, in six stages. The masks come from
PRISM's segmenter or are those supplied with the dataset; the results in this paper use the supplied masks unless
stated. The masks
give 27 measurements of the tumour and its change ($\mathbf{x}_{\mathrm{img}}$); the treatment course as of the scan
day and the patient's characteristics give the clinical-history inputs ($\mathbf{x}_{\mathrm{hist}}$). Three ridge
logistic regressions score the scan on each input set and on both. The evidence statements are the split-conformal set of the combined score, the agreement of the imaging and
clinical-history scores, and the reliability card. Collation is deterministic code that gathers every score,
measurement and statement into one record, each under a fixed key with its source. A local language model writes the
report from the record alone, and a deterministic checker accepts it or rejects it without repair. Colours mark the models as in every figure; the two
input slices are this patient's last two scans. (d) Models were developed on MU by repeated cross validation grouped
by patient, then fitted once and applied unchanged at LUMIERE and UCSF, under rules fixed before any score was
compared with an external label, except the LUMIERE imaging and reduced clinical-history scores, added afterwards
(Table~\ref{tab:models}).}
\label{fig:1}
\end{figure}

% Drawn by scripts/fig_cohorts.py from the system records and the supplied dates (after Tak et al. 2025, Figure 2).
\begin{figure}[!htbp]
\centering
\figslot{fig_cohorts.pdf}{0.50\textheight}{For MU, LUMIERE and UCSF, the time of each evaluation scan and the interval
since the previous scan, by label.}
\caption{\textbf{When the evaluation scans were taken.} One panel per cohort, headed by its evaluation scans,
patients and share of progression scans. Left, the time of each scan, split by label: at MU-Glioma-Post in weeks
after the end of radiotherapy, with the first 12 weeks shaded; at LUMIERE, which records no radiotherapy dates, in
weeks after surgery; and at UCSF-ALPTDG, which records no end of radiotherapy, in weeks after the first surgery.
Right, weeks since the patient's previous scan. Each column shares one axis across cohorts; scans beyond it are drawn
at the edge and counted. Each point is one scan (filled, progression; open, no progression), and the vertical bar
marks the median. MU shows the 146 scans of the intended use and UCSF the second scan of each grade 4 patient.}
\label{fig:cohorts}
\end{figure}

% Sources: NUMBERS.md 2.1, 3, 4; DETECTION_TASK.md; UCSF_PLAN.md; CLINICIAN_BASELINE_PLAN.md; CHRISTODOULOU_REPLICATION.md;
% SKELETON_2026-11-20.md (timestamp-only check, scratchpad timing_only.py).
\begin{table}[!htbp]
\centering
\caption[The scores in this paper]{\textbf{The scores in this paper.} Every fitted score except the published
radiomics pipeline is a ridge logistic regression fitted within the training data, with unknown IDH or MGMT status coded
with wildtype and unmethylated (coding it as its own category changed no AUC clearly), and missing measurements imputed with the training median and an indicator. Status says whether
a score's own AUC was planned before the analysis or exploratory; the comparisons in Table~\ref{tab:2} were planned
even where a score's AUC on MU was exploratory.}
\label{tab:models}
\footnotesize
\zebrarows
\begin{tabularx}{\linewidth}{@{}R{2.9cm}LR{1.3cm}R{3.4cm}@{}}
\toprule
Score & Inputs & Dates & Status \\
\midrule
Clinical-history model & Age, sex, IDH and MGMT status, days since diagnosis, day of surgery, radiotherapy start, end and dose, fractions, whether chemotherapy had started, the 12-week flag, visit index, biopsy before resection and a previous brain tumour & Yes & Planned \\
Reduced clinical-history model & Age, sex, IDH and MGMT status, days since surgery, a flag for the first weeks after radiotherapy & Partly & Planned at UCSF-ALPTDG, exploratory at LUMIERE \\
Imaging model & 27 measurements of the segmented tumour (volumes, shape and intensity, change since the previous scan, volumes relative to the patient's smallest) & No & Exploratory on MU, planned at UCSF-ALPTDG \\
Images from one scan & The 19 measurements of the current scan & No & Exploratory on MU, planned at UCSF-ALPTDG \\
Combined model & Inputs of both models (externally the reduced clinical-history inputs, volumes and their change, and the interval since a previous scan) & Yes & Exploratory on MU, planned and applied unchanged in its reduced form \\
\addlinespace
Published radiomics pipeline & 256 of 4,892 radiomic features of the current scan ranked within training, class weighted gradient boosting tuned over 400 trials, Platt calibration; the authors' code \citep{christodoulou2026} & No & Planned \\
Time-and-enhancement model & Time window after radiotherapy, log enhancing volume and its change since the previous scan & Yes & Planned \\
Timing alone & Days since the end of radiotherapy or visit number, without a fitted model; or a ridge on these, days since diagnosis and the 12-week flag & Yes & Exploratory \\
Clinical data alone & Age, sex, IDH and MGMT status & No & Planned reference \\
Enhancing volume alone & Log enhancing tumour volume, no fitted model & No & Exploratory reference \\
\bottomrule
\end{tabularx}
\end{table}

% Sources: DETECTION_TASK.md (Plan, B1); UCSF_IMAGING.md; NUMBERS.md 2.1; SEGMENTER_STAGE.md section 2;
% SEGMENTER_BRATS24_PLAN.md; refs mahmoud2025, suter2022, fields2024, hatamizadeh2022, deverdier2024, ferreira2024.
\paragraph*{Segmentation and measurements.} PRISM outlines the tumour with a fixed public network, or uses the masks
supplied with each dataset, and computes 27 measurements of the tumour and its change. The masks supplied with each
dataset (the supplied masks) are automatic masks at MU-Glioma-Post, DeepBraTumIA masks with three labels at LUMIERE
and the masks provided with UCSF-ALPTDG (Table~\ref{tab:1}). The primary results use them, so every volume can be
traced to a public file. PRISM's own segmentation stage holds two networks, each run unchanged with its public weights
and never tuned on our cohorts. The first is a Swin UNETR trained on BraTS 2021 \citep{hatamizadeh2022}
(Figure~\ref{fig:1}b). It takes FLAIR, T1c, T1 and T2 as four channels, each normalised over its non-zero voxels,
slides a window of $128^3$ voxels across the brain with half overlap, and thresholds its three sigmoid outputs (tumour
core, whole tumour and enhancing tumour) at 0.5. The enhancing tumour is taken directly, necrosis is the core less the
enhancing tumour, and FLAIR hyperintensity is the whole tumour less the core. This network was trained on scans before
surgery and does not outline the resection cavity, so it supplies 22 of the 27 measurements. The second is the
first-placed model of the BraTS 2024 challenge on post-treatment adult glioma \citep{deverdier2024,ferreira2024}, an
ensemble of nnU-Net and MedNeXt networks that labels the resection cavity as well and so supplies all 27. LUMIERE's
images are stored on an atlas grid with a different axis orientation, so before either network they are reoriented to
the orientation of the training data, and the label map is mapped back. This step is the same for every LUMIERE scan.

From the masks and images, PRISM computes 27 measurements. Nineteen describe the current scan: the log volumes of the
enhancing tumour, necrosis, FLAIR hyperintensity, resection cavity, tumour core and whole lesion, and 13 measures of
shape and of T1 contrast and FLAIR intensity relative to normal appearing brain (Supplementary Table~S7). Eight
describe change since the previous scan and since the patient's smallest volume. Five measurements depend on the
cavity (its volume, whether there is one, the share of enhancing tumour next to it, the distance of the enhancing
tumour from it, and its change since the previous scan).

PRISM also takes the treatment course as of the scan day: days since diagnosis and surgery, radiotherapy start, end
and dose, and whether chemotherapy had started. It places the scan in a time window after radiotherapy (before its
end, within 12 weeks after its end, later, or unknown when dates are missing), and it takes age, sex, and IDH mutation
and MGMT promoter methylation status. Every input is taken as of the scan day. IDH and MGMT status count as unknown at
scans before the diagnosis date (8 grade 4 scans, used only in training). The duration of radiotherapy and the time
since its end enter only after radiotherapy has ended, and the prescribed dose enters from its start.

% Sources: prism/models/formulations.py; prism/experiments/c3_radiomics.py (BASE_PRE); NUMBERS.md 2.1;
% DETECTION_TASK.md (External test); SYSTEM_RESULTS.md (Deviations item 1).
